% v21: current frontier and reusable next steps; stable support identifiers.
\subsection{残る五型と，再開する具体的な入口}
\label{ns20:support-agenda}

現在の高分母 support の必要候補は\emph{五型・六十六標識}である．
再開時は最初に \nolinkurl{00_CONTINUATION_POLICY_ja.txt} を読む．
以下の順位は，面白さとインパクトを最大化するために柔軟に変更する．
全候補は五頂点を使い，最小次数二以上，二点連結，辺数七以上，$b_1\ge3$ である．
各二価頂点の二つの高近傍はともに次数四である．
二価頂点数は零・一・三に限られ，一個の場合は $G_{5,15}$，
三個の場合および七高辺の場合は book $G_{5,13}$ となる．
表の $ij$ は無向辺，$N$ は標識付き個数であり，補 graph の孤立頂点は省略する．
\begin{center}
\footnotesize
\setlength{\tabcolsep}{4pt}
\begin{tabular}{crclcr}
型&$|E|$&次数列&$E$&低辺の補 graph&$N$\\\hline
$G_{5,13}$&7&$(44222)$&$\{01,02,03,04,12,13,14\}$&$K_3$&10\\
$G_{5,14}$&8&$(43333)$&$K_5-\{14,23\}$&$2K_2$&15\\
$G_{5,15}$&8&$(44332)$&$K_5-\{24,34\}$&$P_3$&30\\
$G_{5,16}$&9&$(44433)$&$K_5-\{34\}$&$K_2$&10\\
$G_{5,17}$&10&$(44444)$&$K_5$&$\varnothing$&1
\end{tabular}
\end{center}
全 $1024$ 標識付き graph の現行検算は
\nolinkurl{verify_support_agenda_v20.py} と
\nolinkurl{support_agenda_verification_v20.json} に保持する．
候補数は整数接着，atlas，幾何的非同型類の存在数ではない．

\paragraph{第一候補：book の残余二領域．}
二価 cores $2,3,4$ の低三角形には，負枝がちょうど一つ必要である．
唯一の負枝を $\sigma_N=+$ と正規化すると，残る型は
\[
\begin{array}{c|c|c}
 & \text{node 周期} & (L_{P_1},L_{P_2})\\ \hline
\mathrm I&N^+\to P_1^+\to P_2^+\to N^+&(3,3)\\
\mathrm {II}&N^+\to P_1^+\to P_2^-\to N^+&(4,3)
\end{array}
\]
であり，双方 $\tau=+$，$w\ge135,H\ge w$ である．
全ての二価 cores で $q_{i0}<q_{i1}$ も必要となる．
型IIでは $S_{P_2}=w-1,T_{P_2}=w+1$ と最後の node $(0,0,5)$ が固定される．
この固定差と分母順序を，全 scale の因子式と有限 node offset に同時に課す．
新しい成長不等式と二つの共鳴族の扱いを保持し，単なる探索上限の拡張は優先しない．
二領域の存在・不存在はともに未確定である．
低三角形が残った場合に四高中心 $0,1$ と $q_{01}$ を復元する．

\paragraph{並行候補：$G_{5,15}$ の五つの固定低データ問題．}
この還元では標識を $H=K_5-\{23,24\}$ と取り直し，二価 core を $2$，
共有高分母を $Q=q_{34}$ とする．表の代表とは標識置換だけが異なる．
正例外型は全排除され，通常型 $v=-u$ は四中心・十四配置へ限られる．
負例外型は $Q=3,5$ の二 template となる．
通常型は各固定 $Q$ について cores $2,3,4$ と相互三辺を完全に有限復元できるが，
$Q$ の一様上界を与えるものではない．
次は正 reverse 側の固定低 row と要求 $D_0$ の五問題
\[
((-1,-4,4),5),\quad ((-1,-1,1),-9),\quad
((-1,-2,2),2),\quad ((-1,-2,2),-7),\quad ((-1,-2,2),9)
\]
を無上界で分類する．外部二標識の交換は許している．
負例外では $Q$ が小さくても二高中心の自由度が残るため，別途全尺度を扱う．
$G_{5,14}$ と $G_{5,16},G_{5,17}$ も必要候補として保持する．
四高側へ三高 matching 定理を無条件に拡張しない．

\paragraph{三高 $D$ 問題：単一合同式と残余式の交わり．}
\ref{nhd21:reconstruction} の五データ復元を用いる．
整数尺度には $m=F^2+F+1\mid FE-J$ が必要十分であり，
$N=E^2+EJ+J^2>0$ と $m\mid N$ により各固定 $(s,t,x,y,z)$ の $b$ は有限となる．
ただしノルム条件だけでは合同式に足りず，尺度を復元しても元の pair 残余式を残す．
$j=t-y$ を正と仮定してはいけない．
五変数自体への一様上界，または残余式との無上界矛盾を目標とする．
正整数候補が得られたときは明示的な正 weights まで復元する．

\paragraph{再利用と，解決済みの型．}
$K_4$，二葉三角形，$G_{5,5},G_{5,7},G_{5,9},G_{5,10},G_{5,11},G_{5,12}$ は再探索しない．
reverse data の解析的上界，点・全半直線表示，分母一致と $D$ の奇偶障害は再利用できる．
ただし相手の高次数と例外接着の仮定を毎回確認する．
$G_{5,15},G_{5,16},G_{5,17}$ が $K_4$ を部分 graph に含むだけでは，
support 全体が $K_4$ の場合の排除を移せない．
高分母 support $H$ と素数 support $G_\ell$ を区別する．
dense mod 五の二模型では少なくとも二辺の五進付値一を保つ．
整数候補が残れば actual weights，basin areas，integer framings，作用切片へ進む．
幾何的同値，別退化，局所置換への転換も，研究全体の価値に応じて判断する．
